\documentclass[10pt,a4paper]{amsart}
\usepackage{amsmath,amssymb,amsthm,graphicx}
% Retained minimal preamble for the Stacks Project chapter more-algebra.tex (More on Algebra).
% The chapter file pinned-original-source.tex carries only \input{preamble} on its first line, so no
% preamble text lives inside the chapter itself.  The Stacks Project uses one single shared preamble for
% every chapter; the snapshot stages it identically as sources/stacks-chow/preamble.tex and
% sources/stacks-etale-cohomology/preamble.tex.  The block below is copied mechanically, line for line,
% from sources/stacks-chow/preamble.tex lines 14-18 (the Xy-pic setup) and 154-217 (the native theorem
% declarations together with the mathematical macros).  Nothing was added, removed or redefined; the
% wrapper class and the amsmath/amssymb/amsthm/graphicx base packages are supplied by the assembler.
% Chapter: more-algebra.tex -- More on Algebra.
% For commutative diagrams we use Xy-pic
\usepackage[all]{xy}

% We use 2cell for 2-commutative diagrams.
\xyoption{2cell}

% Theorem environments.
%
\theoremstyle{plain}
\newtheorem{theorem}[subsection]{Theorem}
\newtheorem{proposition}[subsection]{Proposition}
\newtheorem{lemma}[subsection]{Lemma}

\theoremstyle{definition}
\newtheorem{definition}[subsection]{Definition}
\newtheorem{example}[subsection]{Example}
\newtheorem{exercise}[subsection]{Exercise}
\newtheorem{situation}[subsection]{Situation}

\theoremstyle{remark}
\newtheorem{remark}[subsection]{Remark}
\newtheorem{remarks}[subsection]{Remarks}

\numberwithin{equation}{subsection}

% Macros
%
\def\lim{\mathop{\mathrm{lim}}\nolimits}
\def\colim{\mathop{\mathrm{colim}}\nolimits}
\def\Spec{\mathop{\mathrm{Spec}}}
\def\Hom{\mathop{\mathrm{Hom}}\nolimits}
\def\Ext{\mathop{\mathrm{Ext}}\nolimits}
\def\SheafHom{\mathop{\mathcal{H}\!\mathit{om}}\nolimits}
\def\SheafExt{\mathop{\mathcal{E}\!\mathit{xt}}\nolimits}
\def\Sch{\mathit{Sch}}
\def\Mor{\mathop{\mathrm{Mor}}\nolimits}
\def\Ob{\mathop{\mathrm{Ob}}\nolimits}
\def\Sh{\mathop{\mathit{Sh}}\nolimits}
\def\NL{\mathop{N\!L}\nolimits}
\def\CH{\mathop{\mathrm{CH}}\nolimits}
\def\proetale{{pro\text{-}\acute{e}tale}}
\def\etale{{\acute{e}tale}}
\def\QCoh{\mathit{QCoh}}
\def\Ker{\mathop{\mathrm{Ker}}}
\def\Im{\mathop{\mathrm{Im}}}
\def\Coker{\mathop{\mathrm{Coker}}}
\def\Coim{\mathop{\mathrm{Coim}}}

% Boxtimes
%
\DeclareMathSymbol{\boxtimes}{\mathbin}{AMSa}{"02}

%
% Macros for moduli stacks/spaces
%
\def\QCohstack{\mathcal{QC}\!\mathit{oh}}
\def\Cohstack{\mathcal{C}\!\mathit{oh}}
\def\Spacesstack{\mathcal{S}\!\mathit{paces}}
\def\Quotfunctor{\mathrm{Quot}}
\def\Hilbfunctor{\mathrm{Hilb}}
\def\Curvesstack{\mathcal{C}\!\mathit{urves}}
\def\Polarizedstack{\mathcal{P}\!\mathit{olarized}}
\def\Complexesstack{\mathcal{C}\!\mathit{omplexes}}
% \Pic is the operator that assigns to X its picard group, usage \Pic(X)
% \Picardstack_{X/B} denotes the Picard stack of X over B
% \Picardfunctor_{X/B} denotes the Picard functor of X over B
\def\Pic{\mathop{\mathrm{Pic}}\nolimits}
\def\Picardstack{\mathcal{P}\!\mathit{ic}}
\def\Picardfunctor{\mathrm{Pic}}
\def\Deformationcategory{\mathcal{D}\!\mathit{ef}}

% --- appended 2026-10-12 (Chief): the project own definitions of the reference/slogan/history
% display environments, copied mechanically from the shared Stacks preamble lines 8-12
% (\usepackage{verbatim} plus the three \newenvironment lines, which the chapter preamble keeps
% verbatim from the project preamble). Nothing added, removed or redefined: these environments
% carry no mathematical content and the project itself typesets them as comments. Without these
% three lines a counted statement whose own text uses them fails with Environment slogan undefined.
\usepackage{verbatim}
\newenvironment{reference}{\comment}{\endcomment}
\newenvironment{slogan}{\comment}{\endcomment}
\newenvironment{history}{\comment}{\endcomment}
\usepackage{cleveref}
\crefname{originallemma}{original lemma}{original lemmas}
\Crefname{originallemma}{Original lemma}{Original lemmas}
\makeatletter
\newlabel{algebra-lemma-completion-finite-extension}{{07N9}{source}{}{sourcepointer-1}{}}
\newlabel{algebra-lemma-completion-finite-extension@cref}{{[originallemma][0][]07N9}{[1][0][]source}{}{}{}}
\makeatother
\crefname{originallemma}{original lemma}{original lemmas}
\Crefname{originallemma}{Original lemma}{Original lemmas}
\makeatletter
\newlabel{algebra-lemma-quotient-henselization}{{05WQ}{source}{}{sourcepointer-2}{}}
\newlabel{algebra-lemma-quotient-henselization@cref}{{[originallemma][0][]05WQ}{[1][0][]source}{}{}{}}
\makeatother
\crefname{originallemma}{original lemma}{original lemmas}
\Crefname{originallemma}{Original lemma}{Original lemmas}
\makeatletter
\newlabel{lemma-Nagata-local-ring}{{13413}{source}{}{sourcepointer-3}{}}
\newlabel{lemma-Nagata-local-ring@cref}{{[originallemma][0][]13413}{[1][0][]source}{}{}{}}
\makeatother
\crefname{originallemma}{original lemma}{original lemmas}
\Crefname{originallemma}{Original lemma}{Original lemmas}
\makeatletter
\newlabel{lemma-branches}{{31155}{source}{}{sourcepointer-4}{}}
\newlabel{lemma-branches@cref}{{[originallemma][0][]31155}{[1][0][]source}{}{}{}}
\makeatother
\crefname{originallemma}{original lemma}{original lemmas}
\Crefname{originallemma}{Original lemma}{Original lemmas}
\makeatletter
\newlabel{lemma-check-P-ring-maximal-ideals}{{12954}{source}{}{sourcepointer-5}{}}
\newlabel{lemma-check-P-ring-maximal-ideals@cref}{{[originallemma][0][]12954}{[1][0][]source}{}{}{}}
\makeatother
\crefname{originallemma}{original lemma}{original lemmas}
\Crefname{originallemma}{Original lemma}{Original lemmas}
\makeatletter
\newlabel{lemma-completion-normal-local-ring}{{13485}{source}{}{sourcepointer-6}{}}
\newlabel{lemma-completion-normal-local-ring@cref}{{[originallemma][0][]13485}{[1][0][]source}{}{}{}}
\makeatother
\crefname{originallemma}{original lemma}{original lemmas}
\Crefname{originallemma}{Original lemma}{Original lemmas}
\makeatletter
\newlabel{lemma-formal-fibres-normal}{{13215}{source}{}{sourcepointer-7}{}}
\newlabel{lemma-formal-fibres-normal@cref}{{[originallemma][0][]13215}{[1][0][]source}{}{}{}}
\makeatother
\crefname{originallemma}{original lemma}{original lemmas}
\Crefname{originallemma}{Original lemma}{Original lemmas}
\makeatletter
\newlabel{lemma-henselization-P-ring}{{13155}{source}{}{sourcepointer-8}{}}
\newlabel{lemma-henselization-P-ring@cref}{{[originallemma][0][]13155}{[1][0][]source}{}{}{}}
\makeatother
\crefname{originallemma}{original lemma}{original lemmas}
\Crefname{originallemma}{Original lemma}{Original lemmas}
\makeatletter
\newlabel{lemma-henselization-noetherian}{{10979}{source}{}{sourcepointer-9}{}}
\newlabel{lemma-henselization-noetherian@cref}{{[originallemma][0][]10979}{[1][0][]source}{}{}{}}
\makeatother
\crefname{originalproposition}{original proposition}{original propositions}
\Crefname{originalproposition}{Original proposition}{Original propositions}
\makeatletter
\newlabel{proposition-finite-type-over-P-ring}{{12999}{source}{}{sourcepointer-10}{}}
\newlabel{proposition-finite-type-over-P-ring@cref}{{[originalproposition][0][]12999}{[1][0][]source}{}{}{}}
\makeatother
\title{A Noetherian local ring with geometrically normal formal fibres is Nagata and has the same number of branches as its completion}
\author{The Stacks Project Authors}
\date{stacks-project:more-algebra.tex:local-snapshot-20261012; GNU FDL 1.2 or later}
\begin{document}
\maketitle
\setcounter{subsection}{1}
\noindent\textit{Editorial note.} All delivered mathematical text is the author's own, copied from the pinned Stacks Project chapter file more-algebra.tex; no statement or proof step was rewritten, reordered or removed.\par\medskip
\noindent\textit{Editorial note.} The chapter file pinned-original-source.tex carries only \textbackslash{}input\{preamble\} on its first line, so no preamble text lives inside the chapter. The retained file sources/stacks-more-algebra/retained-preamble.tex is a mechanical, line-for-line copy of the project's single shared preamble as staged in the snapshot, together with the project's own definitions of the reference/slogan/history display environments, which the project itself typesets as comments; nothing was added, removed or redefined by this card and no package was added.\par\medskip
\noindent\textit{Editorial note.} The source URL is the source-level Stacks Project URL, as required when more than one card is produced from one source; the per-claim tag is carried by claim\_key.\par\medskip
\noindent\textit{Editorial note.} The counted statement span includes the author's own \textbackslash{}begin\{reference\} block citing the theorem in the literature; it is delivered verbatim and produces no unresolved citation. The author writes 'small detail omitted' inside the counted proof span, kept verbatim. The eleven references inside the counted proof span are recorded as non-load-bearing navigation pointers.\par\medskip
\noindent\textit{Allowed standard prerequisites (editorial).}\begin{itemize}
\item Nagata rings and finiteness of the integral closure in the fraction field. used at proof lines 32242-32243 and 32263-32264. Supplies both the Nagata conclusion and the finiteness that makes the branch count of A equal to the number of maximal ideals of A'.
\item P-rings for a property of formal fibres, and the criterion on maximal ideals. invoked at proof lines 32244-32248. Propagates geometric normality of the formal fibres to the auxiliary rings used in the reduction.
\item Normality of the completion of a local ring with geometrically normal formal fibres. invoked at proof lines 32272-32275. Is what makes the minimal primes of A' tensor\_A A-hat match the factors of the product decomposition.
\item Flatness of the completion map and commutation of completion with finite extensions. used at proof lines 32268-32271 and 32278-32281. Gives the decomposition of A'-hat and the chain of inclusions whose minimal primes are compared.
\end{itemize}\medskip
\ifdefined\hypertarget\hypertarget{sourcepointer-1}{}\fi
\noindent\textit{Source reference.} Original lemma 07N9: focused-seeds/AGP/sources/stacks-snapshot-20261011/tags.txt, source lines 9754--9754. Cross-chapter navigation pointer into the Stacks Project; the standard fact it supplies is applied inside the counted proof span and the argument resting on it stays there.\par\medskip
\ifdefined\hypertarget\hypertarget{sourcepointer-2}{}\fi
\noindent\textit{Source reference.} Original lemma 05WQ: focused-seeds/AGP/sources/stacks-snapshot-20261011/tags.txt, source lines 7594--7594. Cross-chapter navigation pointer into the Stacks Project; the standard fact it supplies is applied inside the counted proof span and the argument resting on it stays there.\par\medskip
\ifdefined\hypertarget\hypertarget{sourcepointer-3}{}\fi
\noindent\textit{Source reference.} Original lemma 13413: sources/stacks-more-algebra/pinned-original-source.tex, source lines 13413--13413. Same-chapter navigation pointer; the statement it names is standard for the section and the argument that consumes it stays inside the counted proof span.\par\medskip
\ifdefined\hypertarget\hypertarget{sourcepointer-4}{}\fi
\noindent\textit{Source reference.} Original lemma 31155: sources/stacks-more-algebra/pinned-original-source.tex, source lines 31155--31155. Same-chapter navigation pointer; the statement it names is standard for the section and the argument that consumes it stays inside the counted proof span.\par\medskip
\ifdefined\hypertarget\hypertarget{sourcepointer-5}{}\fi
\noindent\textit{Source reference.} Original lemma 12954: sources/stacks-more-algebra/pinned-original-source.tex, source lines 12954--12954. Same-chapter navigation pointer; the statement it names is standard for the section and the argument that consumes it stays inside the counted proof span.\par\medskip
\ifdefined\hypertarget\hypertarget{sourcepointer-6}{}\fi
\noindent\textit{Source reference.} Original lemma 13485: sources/stacks-more-algebra/pinned-original-source.tex, source lines 13485--13485. Same-chapter navigation pointer; the statement it names is standard for the section and the argument that consumes it stays inside the counted proof span.\par\medskip
\ifdefined\hypertarget\hypertarget{sourcepointer-7}{}\fi
\noindent\textit{Source reference.} Original lemma 13215: sources/stacks-more-algebra/pinned-original-source.tex, source lines 13215--13215. Same-chapter navigation pointer; the statement it names is standard for the section and the argument that consumes it stays inside the counted proof span.\par\medskip
\ifdefined\hypertarget\hypertarget{sourcepointer-8}{}\fi
\noindent\textit{Source reference.} Original lemma 13155: sources/stacks-more-algebra/pinned-original-source.tex, source lines 13155--13155. Same-chapter navigation pointer; the statement it names is standard for the section and the argument that consumes it stays inside the counted proof span.\par\medskip
\ifdefined\hypertarget\hypertarget{sourcepointer-9}{}\fi
\noindent\textit{Source reference.} Original lemma 10979: sources/stacks-more-algebra/pinned-original-source.tex, source lines 10979--10979. Same-chapter navigation pointer; the statement it names is standard for the section and the argument that consumes it stays inside the counted proof span.\par\medskip
\ifdefined\hypertarget\hypertarget{sourcepointer-10}{}\fi
\noindent\textit{Source reference.} Original proposition 12999: sources/stacks-more-algebra/pinned-original-source.tex, source lines 12999--12999. Same-chapter navigation pointer; the statement it names is standard for the section and the argument that consumes it stays inside the counted proof span.\par\medskip
\begin{lemma}
\label{lemma-geometrically-normal-formal-fibres-number-of-branches}
\begin{reference}
\cite[Theorem 2.3]{Beddani}
\end{reference}
Let $(A, \mathfrak m)$ be a Noetherian local ring. If the formal
fibres of $A$ are geometrically normal (for example if $A$ is
excellent or quasi-excellent), then $A$ is Nagata
and the number of (geometric) branches of $A$ and $A^\wedge$ is the same.
\end{lemma}

\begin{proof}
Since a normal ring is reduced, we see that $A$ is Nagata by
Lemma \ref{lemma-Nagata-local-ring}. In the rest of the proof
we will use Lemma \ref{lemma-formal-fibres-normal},
Proposition \ref{proposition-finite-type-over-P-ring}, and
Lemma \ref{lemma-check-P-ring-maximal-ideals}. This tells us
that $A$ is a P-ring where $P(k \to R) = $``$R$ is geometrically
normal over $k$'' and the same is true for any (essentially of) finite type
$A$-algebra.

\medskip\noindent
Let $\mathfrak q \subset A$ be a minimal prime. Then
$A^\wedge/\mathfrak q A^\wedge = (A/\mathfrak q)^\wedge$
and $A^h/\mathfrak qA^h = (A/\mathfrak q)^h$
(Algebra, Lemma \ref{algebra-lemma-quotient-henselization}).
Hence the number of branches of $A$ is the sum of the
number of branches of the rings $A/\mathfrak q$ and
similarly for $A^\wedge$. In this way we reduce
to the case that $A$ is a domain.

\medskip\noindent
Assume $A$ is a domain. Let $A'$ be the integral closure of $A$
in the fraction field $K$ of $A$. Since $A$ is Nagata, we see that
$A \to A'$ is finite. Recall that the number of branches of
$A$ is the number of maximal ideals $\mathfrak m'$ of $A'$
(Lemma \ref{lemma-branches}). Also, recall that
$$
(A')^\wedge = A' \otimes_A A^\wedge =
\prod\nolimits_{\mathfrak m' \subset A'} (A'_{\mathfrak m'})^\wedge
$$
by Algebra, Lemma \ref{algebra-lemma-completion-finite-extension}.
Because $A'_{\mathfrak m'}$ is a local ring whose formal
fibres are geometrically normal, we see that
$(A'_{\mathfrak m'})^\wedge$ is normal
(Lemma \ref{lemma-completion-normal-local-ring}).
Hence the minimal primes of $A' \otimes_A A^\wedge$
are in $1$-to-$1$ correspondence with the factors in the
decomposition above. By flatness of $A \to A^\wedge$ we have
$$
A^\wedge \subset A' \otimes_A A^\wedge \subset K \otimes_A A^\wedge
$$
Since the left and the right ring have the same set of minimal
primes, the same is true for the ring in the middle (small
detail omitted) and this finishes the proof.

\medskip\noindent
To see this is true for the number of geometric branches
we use the result for branches, the fact that the formal
fibres of $A^{sh}$ are geometrically normal
(Lemmas \ref{lemma-formal-fibres-normal} and
\ref{lemma-henselization-P-ring})
and the fact that $(A^{sh})^\wedge = ((A^\wedge)^{sh})^\wedge$
by Lemma \ref{lemma-henselization-noetherian}.
\end{proof}

\end{document}
